\section{Trí tuệ nhân tạo}

\subsection{Mô hình FitDiT}
FitDiT (Fit Diffusion Transformer) là mô hình học sâu chuyên biệt 
cho bài toán Virtual Try-On, được phát triển bởi nhóm nghiên cứu 
từ Tencent và Đại học Fudan \cite{fitdit_paper}. Mô hình sử dụng 
kiến trúc Diffusion Transformer (DiT) thay thế cho U-Net truyền 
thống, cho phép tập trung xử lý nhiều hơn ở các đặc trưng độ 
phân giải cao, giải quyết hai thách thức chính là bảo toàn kết 
cấu vải (texture-aware maintenance) và khớp kích thước trang phục 
(size-aware fitting). Trong đề tài, FitDiT nhận đầu vào là ảnh 
người dùng và ảnh quần áo, sau đó sinh ra ảnh kết quả thể hiện 
người dùng đang mặc trang phục đó một cách chân thực.

\subsection{RAG (Retrieval-Augmented Generation)}
RAG là kỹ thuật kết hợp giữa truy xuất thông tin (retrieval) và 
sinh văn bản (generation), giúp mô hình ngôn ngữ trả lời dựa trên 
dữ liệu thực tế được cung cấp thay vì chỉ dựa vào kiến thức được 
huấn luyện sẵn \cite{rag_lewis_2020}. Trong đề tài, RAG được ứng 
dụng cho chatbot tư vấn thời trang, sử dụng pgvector để lưu trữ 
và tìm kiếm vector embeddings của thông tin sản phẩm, từ đó cung 
cấp câu trả lời chính xác và phù hợp với danh mục hàng hóa thực 
tế của hệ thống.

\subsection{Ollama}
Ollama là công cụ mã nguồn mở cho phép triển khai và chạy các 
mô hình ngôn ngữ lớn (LLM) trên máy cục bộ mà không cần kết nối 
Internet hay dịch vụ đám mây \cite{ollama_docs}. Trong đề tài, 
Ollama được sử dụng để chạy mô hình Qwen2.5 7B trên máy chủ nội 
bộ, phục vụ cho chức năng chatbot tư vấn thời trang, đảm bảo tính 
bảo mật dữ liệu người dùng và giảm chi phí vận hành.

\subsection{pgvector}
pgvector là extension mã nguồn mở cho PostgreSQL, bổ sung khả năng 
lưu trữ và tìm kiếm tương đồng vector (vector similarity search) 
trực tiếp trong cơ sở dữ liệu quan hệ \cite{pgvector_github}. 
Trong đề tài, pgvector được tích hợp vào Neon PostgreSQL để lưu 
trữ vector embeddings của sản phẩm (1536 chiều), phục vụ cho 
việc tìm kiếm ngữ nghĩa trong luồng RAG của chatbot tư vấn 
thời trang.